\documentclass[twoside,a4paper]{article}
\usepackage{amsmath}
\usepackage{amsfonts}
%\usepackage{mathabx} %for the \wdcheck
\usepackage{amssymb}
\usepackage[labelfont=bf]{caption}
\usepackage[english]{babel}
\usepackage[utf8]{inputenc}
\usepackage[OT1]{fontenc}
\usepackage{graphicx}
\usepackage{latexsym}
\usepackage{subcaption} 
\usepackage[inner=36pt, textwidth=345pt]{geometry}
\usepackage{subfiles}
\usepackage{siunitx}
\usepackage{tikz}
\usetikzlibrary{calc}
\usetikzlibrary{scopes}
\usepackage{ifthen} % for the tikz loops
\usepackage[toc,page]{appendix}
\usetikzlibrary{shapes.geometric}
\usetikzlibrary{decorations.pathmorphing}
\usepackage{hyperref}
\usepackage[makeroom]{cancel}
\usepackage{pdfpages}
\usepackage{bm}
\usepackage{layout}
\usepackage{ relsize }
\usepackage{float}
\usepackage{cancel}
\usepackage{stmaryrd}
% \usepackage{TamwynLatexLib}
\usepackage{fancyhdr}
\usepackage[explicit]{titlesec}   
\usepackage{mathtools}
\usepackage{pgfplots}
\pgfplotsset{width=10cm,compat=1.9}
\usepackage{pdfpages}
\usetikzlibrary{shadows}
\usepackage{bbm} %for the \mathbb{1}
\newcommand{\indc}{\mathlarger{\mathlarger{\mathbbm{1}\normalsize}}}

% \usepackage[style=numeric, sorting=none]{biblatex} %backend=biber, sortlocale=en_EN

%\input{C:/Users/Tamwyn/Documents/CustomLatexTextMF/mytexmf/tex/latex/UniLaTeXPackage/Arto/ArtoMain.tex}


\geometry{ 
    bindingoffset=1.0471975512cm, 
    left=2.9cm, 
    right=2.6cm, 
    top=1.6cm, 
    bottom=2.5cm
}

\renewcommand{\footnoterule}{%
    \vspace*{-3pt} % Vertical space before the line
    \begin{tikzpicture}
        \coordinate (A) at (0.2\textwidth,-0.5);
        \coordinate (B) at (0,-0.5);
        \draw[black, line width = 0.5pt] (B) -- (A);
        \draw ($(A) + (0.1,0 )$) rectangle ++(0.1,0.1);
        \draw ($(B) - (0.3,0 )$) rectangle ++(0.2,0.2);
    \end{tikzpicture}

    \vspace*{3pt} % Vertical space after the line
}
%\thispagestyle{empty}
% \setlength\parindent{0pt}
\pagestyle{fancy}
\fancyhf{} % Clear default header and footer
% Customize footer with design around the page number
\fancyfoot[C]{
    \begin{tikzpicture}[baseline=(X.base)]
        \node (X) at (0,0) [] {\thepage};
        \coordinate (A) at ($(X.north west) + (-0.1,-0.2)$);
        \coordinate (A2) at ($(X.north west) + (-0.5,-0.2)$);
        \coordinate (A3) at ($(X.north west) + (-0.6,-0.2)$);
        \coordinate (B) at ($(X.north east) + (0.1,-0.2)$);
        \coordinate (B2) at ($(X.north east) + (+0.5, -0.2)$);
        \coordinate (B3) at ($(X.north east) + (+0.6, -0.2)$);
        \coordinate (C) at ($(X.south east) + (0.1,0.15)$);
        \coordinate (C2) at ($(X.south east) + (0.2,0.15)$);
        \coordinate (D) at ($(X.south west) + (-0.1,0.15)$);
        \coordinate (D2) at ($(X.south west) + (-0.2,0.15)$);

        \draw (A) -- (A2);
        \draw (B) -- (B2);
        \draw (C) -- (C2);
        \draw (D) -- (D2);
        \draw (A3) rectangle ++(-0.1,0.1);
        \draw (B3) rectangle ++(0.1,0.1);
    \end{tikzpicture}
}
\renewcommand{\headrulewidth}{0pt}
\DeclareSIUnit{\site}{\text{site}}
% Define a custom macro for the equation with a condition

% \pgfmathdeclarefunction{ifthenelsefpu}{3}{%
%   \pgfmathparse{#1*#2 + !#1*#3}%
% }

\pgfmathdeclarefunction{ifthenelsefpu}{3}{%
  \pgfmathparse{#1 == 1 ? #2 : #3}%
}
% \pgfmathdeclarefunction{f}{3}{%
% \pgfmathparse{ifthenelse(abs(cos(#1-pi) + cos(#2-pi) - (#3/2)) < 1,  #3, nan)}%
% }
\pgfmathdeclarefunction{f}{3}{%
  \pgfmathparse{ifthenelse(abs(cos(deg(#1-pi)) + cos(deg(#2-pi)) - (#3/2)) < 5e-2,  #3/2 , nan)}%
}
% Define the function fd
\pgfmathdeclarefunction{fd}{1}{%
  \pgfmathparse{ifthenelse(#1 < 1e-3, #1/2, nan)}%
}


% \pgfmathdeclarefunction{g}{1}{%
%   \pgflibraryfpuifactive{%
%     \pgfkeys{/pgf/fpu=false}%
%     \pgfmathfloattofixed{#1}%
%     \let\x=\pgfmathresult%
%   }%
%   {%
%     \pgfmathparse{#1}%
%     \let\x=\pgfmathresult%
%   }%
%   \pgfmathparse{ifthenelse(\x<0,(\x)^2,\x)}%
% }

% Optional: Adjust spacing before/after sections
\titlespacing*{\section}{0pt}{1em}{1em}

% Temporarily disable custom formatting when including subfiles


\definecolor{TamGreen}{HTML}{00C4F6}
\definecolor{TamTangerine}{HTML}{325BCB}
\definecolor{TamDarkBlue}{HTML}{666A86}
\definecolor{TamDarkBlue2}{HTML}{788AA3}
\definecolor{TamLightGreen}{HTML}{92B6B1}
\definecolor{TamLightGY}{HTML}{B2C9AB}
\definecolor{TamYellow}{HTML}{E8DDB5}
\definecolor{Brillou1}{HTML}{2D70B3}
\definecolor{Brillou2}{HTML}{3D8F4A}
\definecolor{Brillou3}{HTML}{FA7E19}
\definecolor{Brillou4}{HTML}{C74440}

\newcommand{\rem}[1]{\textcolor{TamGreen}{#1}}
\newcommand{\im}{\large\mathfrak{i}\hspace*{1.3pt}\normalsize}
\newcommand{\dd}{\mathrm{d}}